\documentclass[../../main.tex]{subfiles}

\begin{document}

\section{Group theory}
\label{sec:group_theory}

\subsection{Basic definitions}
\label{subsec:gt_basic_definitions}

\begin{definition}
    A \emph{representation} of dimension $n$ of the abstract group $G$ is defined as a homomorphism $D: G \rightarrow \textnormal{GL}(n, \mathbb{C})$, the group of non-singular $n \times n$ matrices with complex entries. 

    In other words it is a mapping under which $g \mapsto D(g)$, perserving the group structure:
    \begin{equation}
        D(g_1 g_2) = D(g_1) D(g_2) \qquad \text{(homomorphism)}.
    \end{equation}
    The mapping is necessarily into the set of non-singular matrices, since each matrix must be invertible: $D(g^{-1}) = (D(g))^{-1}$.
\end{definition}

\begin{definition}
    Two elements $a$ and $b$ of a group $G$ are \emph{conjugate} if there exists an element $g \in G$ such that $a=gbg^{-1}$. The element $g$ is called the \emph{conjugating element}. This is an example of an \emph{equivalence relation}\footnote{A relation must be \emph{reflexive}, \emph{symmetric} and \emph{transitive} to be an equivalence relation.}.

    \emph{Conjugacy classes} are defined as
    \begin{equation}
        (a) = \{b \mid \exists g \in G: b=gag^{-1} \}
    \end{equation}
    and are just one example of \emph{equivalence classes}. The equivalence classes are disjoint sets, that cover the whole set $G$.
\end{definition}

\begin{definition}
    The character of a representation $D$ of a group $G$ is the set $\chi = \{\chi(g) \mid g \in G\}$, where $\chi(g)$ is the trace of the representation matrix $D(g)$:
    \begin{equation}
        \chi(g) = \textnormal{Tr}(D(g)). 
    \end{equation}
\end{definition}
Characters are useful, since they are invariant under equivalence transformations $D(g) \rightarrow S D(g) S^{-1}$ due to the cyclic property of the trace.

\begin{definition}
    A representation of dimension $n+m$ is said to be \emph{reducible} if $D(g)$ takes the form
    \begin{equation}
        D(g) = \begin{pmatrix}
            A(g) & C(g) \\
            O & B(g)
        \end{pmatrix} \label{eq:gt_reducible_representation}
    \end{equation} 
    $\forall g \in G$, where $A, C$ and $B$ are submatrices of dimension $m \times m$, $m \times n$ and $n \times n$ respectively, and $O$ denotes a null matrix of dimension $n \times m$.
\end{definition}
From $D(g)D(g') = D(gg')$ follows
\begin{align*}
    A(gg') &= A(g)A(g') \\
    B(gg') &= B(g)B(g') ,
\end{align*} 
so that $\{A(g)\}$ and $\{B(g)\}$ constitute $[m]$ and $[n]$ representations of $G$ respectively. For finite groups it can be shown (see Maschke's theorem, below) that under equivalence, $C$ can be taken to be a null matrix. The representation is then said to be \emph{completely reducible} or \emph{decomposable}, and we write
\begin{equation*}
    D(g) = A(g) \oplus B(g) .
\end{equation*}
It may be that the representations $A$ and $B$ are themselves decomposable, in which case it is natural to continue the process, which will terminate when we reach the level of \emph{irreducible representations} (irrep), namely representations which cannot be further reduced. Thus, the matrix representation $D$ of the abstract group $G$ can be written as a direct sum of irreducible representations.

\begin{definition}
    A map $T: V \rightarrow V$ is \emph{linear} if
    \begin{equation}
        T(\alpha \vec u + \beta \vec v) = \alpha(T \vec u) + \beta (T \vec v)
    \end{equation}
    for all $\vec u, \vec v \in V; \alpha, \beta \in \mathbb{C}$.
\end{definition}

\paragraph{G-module} 
It turns out to be extremely useful to think of a representation of a group $G$ as a homomorphism, mapping $G$ into the group of non-singular\footnote{Non-singular means, that $T(\vec v) = 0 \Rightarrow \vec v = \vec 0$, which is equivalent to $T$ being invertible.} linear transformations of a vector space $V$. That is, every element $g$ of $G$ is mapped into a linear transformation $T(g)$ with the all-important property that 
\begin{equation}
    T(g'g) = T(g')T(g) \qquad \text{(homomorphism)}.
\end{equation}
A vector space in which a group $G$ acts in this way is called a \emph{G-module}.

\paragraph{Reducibility}
In this language reducibility corresponds to the existence of a \emph{submodule}, i.e. a subspace $U$ of the vector space $V$ which itself is closed under the action of the group:
\begin{equation}
    \vec u \in U \subset V \Rightarrow T(g)\vec u \in U .
\end{equation}
Using some basis of $V$ to express $T(g)$ as a matrix $D(g)$, one can easily show that the aforementioned closure of $U$ is equivalent to $D(g)$ having the form as in \cref{eq:gt_reducible_representation}.

\paragraph{Unitary transformations}
A linear transformation $T$ acting in the vector space $V$ is said to be \emph{unitary} if
\begin{equation}
    (T\vec u, T \vec v) = (\vec u, \vec v) \qquad \forall\vec u, \vec v \in V ,
\end{equation}
where $(\cdot\, , \cdot)$ denotes the scalar product. This is equivalent to
\begin{equation}
    (T\vec u, \vec v) = (\vec u, T^{-1} \vec v) , \label{eq:gt_def_unitary_trafo}
\end{equation} 
since $T$ is invertible. 

\paragraph{Complete reducibility}
Again, let $U$ be a submodule of the $G$-module $V$, closed under the action of the group: $\vec u \in U \subset V \Rightarrow T(g)\vec u \in U$. Use the scalar product to obtain an orthonormal basis of $U$ and of $W = V \setminus U$ such that W is the orthogonal complement of $U$:
\begin{equation}
    W = \{\vec w \in V \mid (\vec w, \vec u) = 0 \quad \forall \vec u \in U \}.
\end{equation}
In coordinate-free language, complete reducibility corresponds to the closure of $W$ as well as $U$ under the aciton of the group. This will always be the case if the $T(g)$ are unitary with respect to the scalar product we are using. For then, from \cref{eq:gt_def_unitary_trafo},
\begin{equation*}
    (T(g)\vec w, \vec u) = (\vec w, T^{-1}(g) \vec u) = (\vec w, T(g^{-1}) \vec u) = (\vec w, \vec u') = 0 ,
\end{equation*}
which shows that $W$ is closed under the action of the group, as required. This translates to the submatrix $C(g)$ being the null matrix.

\paragraph{Maschke's theorem}
All reducible representations of a finite group are completely reducible (decomposable). Since we already know, that all \emph{unitary} representations are decomposable, one only has to show that for a finite group it is always possible to choose a scalar product which makes the representation unitary. This then implies that every reducible matrix representation $D$ is equivalent to a reducible unitary representation $D'$.

\paragraph{Schur's first lemma} \label{para:gt_schurs_first_lemma}
In matrix form this states that any matrix which commutes with all the matrices of an irreducible representation must be a multiple of the unit matrix, i.e.
\begin{equation}
    B D(g) = D(g) B \quad \forall g \in G \Rightarrow B = \lambda \mathds{1},
\end{equation}
with $\lambda \in \mathbb{C}$. As a statement about linear operators it says that 
\begin{equation}
    \hat B T(g) = T(g) \hat B \quad \forall g \in G \Rightarrow \hat B = \lambda E ,
\end{equation}
where $\hat B: U \rightarrow U$ is the linear operator corresponding to B, and E is the identity map $\vec u \mapsto \vec u$.

\paragraph{Schur's second lemma}
This concerns two inequivalent irreducible representations $D, D'$, and in matrix form states that
\begin{equation}
    B D(g) = D'(g)B \quad \forall g \in G \Rightarrow B = 0 . \label{eq:gt_schurs_second_lemma_matrix}
\end{equation} 
The corresponding linear operator $\hat B$ will in this case be a map $\hat B: U \rightarrow U'$ between the two vector spaces $U$ and $U'$ supporting the two representations, and the coordinate-free version of \cref{eq:gt_schurs_second_lemma_matrix} is
\begin{equation}
    \hat B T(g) = T'(g) \hat B \quad \forall g \in G \Rightarrow \hat B = \hat O ,
\end{equation}
where $\hat O$ is the linear operator which maps every vector in $U$ onto the null vector in $U'$: $\hat O \vec u = \vec 0'$ for all $ \vec u \in U$.


\paragraph{The fundamental orthogonality theorem}
Let $U_\nu$ and $U_\mu$ be two $G$-modules carrying inequivalent irreducible representations for a finite group $G$. For the matrix representations then holds
\begin{equation}
    \sum_g D^{(\mu)}_{ir} (g) D^{(\nu)}_{sj} (g^{-1}) = \frac{[g]}{n_\mu} \delta^{\mu \nu} \delta_{ij} \delta_{rs} .
\end{equation}

\paragraph{Orthogonality of characters}
Recall that the character of a representation $D$ is the set $\{\chi(g)\}$, where $\chi(g)$ is the trace of the matrix $D(g)$. These traces have the following properties:
\begin{enumerate}
    \item $\chi$ is the same for equivalent representations, for which $D'(g) = SD(g)S^{-1}$.
    \item $\chi$ is the same for conjugate elements, since $D(hgh^{-1}) = D(h)D(g)(D(h))^{-1}$.
    \item If $D$ is unitary, i.e. $D^{-1} = D^\dagger$, then
    \begin{equation}
        \chi(g^{-1}) = \text{Tr}((D(g))^{-1}) = \text{Tr}(D(g)^\dagger) = \chi^*(g),
    \end{equation}
    which is true for all finite groups.
\end{enumerate}

In terms of characters the orthogonality theorem reads
\begin{equation}
    \frac{1}{[g]} \sum_g \chi^{(\mu)}(g) \chi^{(\nu)^*}(g) = \delta^{\mu \nu}. \label{eq:gt_character_orthogonality}
\end{equation}
Since the characters of conjugate elements are equal, all elements of a conjugacy class will have the same character. Let $\chi_i$ denote the distinct characters where $i=1,\dots,k$ corresponds to the $k$ conjugacy classes $K_i$. \Cref{eq:gt_character_orthogonality} can thus be rewritten as
\begin{equation}
    \frac{1}{[g]} \sum_i k_i \chi_i^{(\mu)}(g) \chi_i^{(\nu)^*}(g) = \delta^{\mu \nu} ,
\end{equation}
with $k_i$ denoting the number of elements in the conjugacy class $K_i$. 

One can prove that the number of inequivalent irreducible representations $r$ is equal to the number of conjugacy classes $k$.


\paragraph{Decomposition of reducible representations}
For a finite group any reducible representation is completely reducible into a direct sum of irreps (Maschke's theorem), which means that the representation matrices can be cast in block diagonal form, the non-zero diagonal blocks being the matrices of irreducible representations. A given irrep may appear more than once in this decomposition: for example, a 5-dimensional representation could decompose into the trivial representation and two copies of the same 2-dimensional irrep. In the general case the decomposition is written
\begin{equation}
    D = \bigoplus_{\nu} a_\nu D^{(\nu)}  ,
\end{equation}
where the non-negative integer $a_\nu$ denotes the number of times the irrep $D^{(\nu)}$ appears in the decomposition. One can show that
\begin{align}
    a_\mu &= \frac{1}{[g]} \sum_g \chi (g) \chi^{(\mu)}(g^{-1}) \\
    &= \frac{1}{[g]} \sum_g \chi (g) \chi^{(\mu)^*}(g) .
\end{align} 

\paragraph{Construction of the character table}
The characters of the irreducible representations of a finite group are conveniently presented in the form of a table, the rows of which correspond to the different irreps, and the columns to the conjugacy classes of the group. In the construction of this table the principal tools we use are:
\begin{enumerate}
    \item number of irreps = number of conjugacy classes: $r=k$;
    \item $\sum_\mu n_\mu^2 = [g]$ ($n_\mu$ is the dimension of the irrep);
    \item orthogonality: $\sum_i k_i \chi_i^{(\mu)} \chi_i^{(\mu)^*}$;
    \item any other information we can come by (!). 
\end{enumerate}
In the last category is the fact that for one-dimensional representations the characters are the same thing as the matrices, and hence the characters themselves must mimic the group multiplication properties.

\paragraph{Abelian groups}


\emph{All irreducible representations of a finite Abelian group $G$ are one-dimensional.}

Take some irreducible matrix representation $D^{(\sigma)}$ of the group $G$\footnote{Each finite group has at least the trivial representation $\rho: G \rightarrow \text{GL}(1, \mathbb{C})$, with $\rho(g) = 1$ for all $g \in G$.}. The matrix $D^{(\sigma)}(g)$ commutes with all the matrices of the representation $D^{(\sigma)}$. Therefore, due to Schur's first lemma, the matrix must be a multiple of the identity: $D^{(\sigma)}(g) = \lambda_g \mathds{1}$. Thus, all the matrices $D^{(\sigma)}(g)$ are diagonal and would therefore be reducible if the dimension is greater than one, which would be a contradiction to our assumption that $D^{(\sigma)}$ is an irrep. 

\subsection{Translation operator in quantum mechanics}
Wiki:

The translation operator $\hat T(\vec x)$ moves particles and fields by the amount $\vec x$ (e.g. $\vec x \in \mathbb R^3$). Therefore, if a particle is in an eigenstate $\ket{r}$ of the position operator, then after $\hat T(\vec x) = \int d\vec r \ket{\vec r + \vec x}\bra{\vec r}$ acts on it, the particle is at the position $\vec r + \vec x$:
\begin{equation}
    \hat T (\vec x) \ket{\vec r} = \ket{\vec r + \vec x}
\end{equation}

The set $\mathcal T$ of translation operators $\hat T(\vec x)$ for all $\vec x$, with the operation of multiplication defined as the result of successive translations, satisfies all the axioms of a group. In order to see this, let $\vec x, \vec x', \vec r \in \mathbb R^3$. Then 
\begin{equation*}
    \hat T (\vec x) \hat T (\vec x') \ket{\vec r} = \hat T (\vec x) \ket{\vec r + \vec x'} = \ket{(\vec r + \vec x') + \vec x} = \hat{T}(\vec x + \vec x') \ket{\vec r},
\end{equation*}
from which directly follows closure under the action of the group. The existence of an identity and of an inverse for each element, and associativity follows directly from the fact, that $\mathbb R^3$ with addition is a group itself. 


% Plan: 
% Abelian group since addition in R^3 commutative -> one dimensional irreps
% Restriction to lattice with pbc -> finite group
% Restriction to one dimension -> cyclic group with generator translation by one lattice constant
% Application to an Hamiltonian: if T^-1 H T = H <=> [H, T] = 0 (since T unitary)
%   -> H and T have common eigenstates -> can be labled by eigenvalues of H or/and (??) T
%   -> H can be written in block diagonal form using irreps of T 
%   -> characters of irreps as good quantum numbers since ??
%   -> why are the crystal/lattice momenta equal to the characters and thus the quantum numbers?
%   -> how do we know which k sector corresponds to state of lowest energy (i.e. the ground state)?
%   -> how does the symmetrizer projector enter here?
%   -> what is the connection to symm_exp_sum?
\end{document}
